%%%%%%%%%%%%%%%%%%%%%%%%%%%%%%%%%%%%%%%%%
% Ecosystem Stewardship — Paper 2
% Based on Paper 1 template (innovation-review)
%%%%%%%%%%%%%%%%%%%%%%%%%%%%%%%%%%%%%%%%%

%----------------------------------------------------------------------------------------
%	PACKAGES AND OTHER DOCUMENT CONFIGURATIONS
%----------------------------------------------------------------------------------------

\documentclass[10pt,twoside,twocolumn,a4paper]{article}

\usepackage[sc]{mathpazo}
\usepackage[T1]{fontenc}
\linespread{1.0}
\usepackage{microtype}

\usepackage[english]{babel}
\usepackage{csquotes}

\usepackage[hmarginratio=1:1,top=25mm,bottom=25mm,left=15mm,right=15mm,columnsep=15pt]{geometry}
\usepackage[hang,small,labelfont=bf,up,textfont=it,up]{caption}

\usepackage{graphicx}
\usepackage{booktabs}
\usepackage{adjustbox}
\usepackage{lettrine}
\usepackage{parskip}
\usepackage{float}

\usepackage{tikz}
\usetikzlibrary{shapes.geometric, arrows.meta, positioning, fit, backgrounds}

\usepackage{enumitem}
\setlist[itemize]{noitemsep}

\usepackage{abstract}
\renewcommand{\abstractnamefont}{\normalfont\bfseries}
\renewcommand{\abstracttextfont}{\normalfont\small\itshape}

\usepackage{titlesec}
\renewcommand\thesection{\Roman{section}}
\titleformat{\section}[block]{\large\scshape\centering}{\thesection.}{1em}{}
\titleformat{\subsection}[block]{\large}{\thesubsection.}{1em}{}

%----------------------------------------------------------------------------------------
%	HEADERS AND FOOTERS
%----------------------------------------------------------------------------------------
\usepackage{fancyhdr}
\pagestyle{fancy}
\fancyhead{}
\fancyfoot{}
\fancyhead[C]{Ecosystem Stewardship}
\fancyfoot[RO,LE]{\thepage}

\usepackage{titling}
\usepackage{hyperref}

%----------------------------------------------------------------------------------------
%	BIBLIOGRAPHY
%----------------------------------------------------------------------------------------
\usepackage[backend=biber,
    style=authoryear,
    minnames=1,
    maxnames=2,
    natbib=true]{biblatex}

\renewcommand*{\bibfont}{\footnotesize}

\addbibresource{references.bib}

\AtEveryBibitem{%
  \clearfield{issn}
  \clearfield{isbn}
  \clearfield{note}
  \clearfield{eprint}
  \ifentrytype{book}{
    \clearfield{pages}
  }{}
  \ifentrytype{misc}{}{%
    \clearfield{url}
  }
}

\def\mytitle{From Vendors to Stewards: How Non-Platform IT Suppliers Orchestrate Client Innovation Ecosystems}

\def\myabstract{\noindent
How do IT suppliers without platform ownership orchestrate client innovation ecosystems? Existing theory offers no answer: transaction cost economics assumes supplier opportunism, dynamic capabilities theory assumes internal control, ecosystem strategy assumes platform-based orchestration, and stewardship theory has not been extended beyond intra-firm relationships. This paper introduces \emph{Ecosystem Stewardship} as a theoretical construct through abductive theory building, using a systematic literature review as evidence of theoretical inadequacy. We identify a four-way structural compression --- hyperscalers, large integrators, SaaS vendors, and generative AI --- that leaves specialist IT firms with one defensible position: contextual intimacy with the client's ecosystem. We propose four mechanisms through which stewardship operates: External Paradox Resolution, Portfolio Transfer, Boundary Management, and Innovation Risk Orchestration. Each mechanism extends existing theory by relaxing a core assumption that blocks explanation of the observed behavior. The framework provides testable propositions and boundary conditions specifying when ecosystem stewardship creates value and when it fails. \\
\\
\textbf{Keywords}: ecosystem stewardship, digital innovation, non-platform orchestration, meta-dynamic capabilities, abductive theory building, contextual intimacy
}

%----------------------------------------------------------------------------------------
%	TITLE SECTION
%----------------------------------------------------------------------------------------

\setlength{\droptitle}{-4\baselineskip}

\pretitle{\begin{center}\Huge \bfseries}
\posttitle{\end{center}}

\title{\mytitle}
\author{%
\textsc{Ilja Heitlager}\thanks{Corresponding author} \\[1ex]
\normalsize Eindhoven University of Technology \\
\normalsize \href{mailto:i.heitlager@tue.nl}{i.heitlager@tue.nl} \\
\textsc{A. Georges L. Romme} \\[1ex]
\normalsize Eindhoven University of Technology \\
\normalsize \href{mailto:a.g.l.romme@tue.nl}{a.g.l.romme@tue.nl}
}
\date{\today}

\renewcommand{\maketitlehookd}{%
\begin{abstract}
\myabstract
\end{abstract}
}

%----------------------------------------------------------------------------------------

\begin{document}

\maketitle

%----------------------------------------------------------------------------------------
%	ARTICLE CONTENTS
%----------------------------------------------------------------------------------------
\section{Introduction}

\subsection{The Structural Compression of IT Services}

The IT services industry is undergoing a structural compression that is eliminating the traditional competitive positions available to specialist firms. This compression operates through four simultaneous forces, each commoditizing a different layer of the value chain.

Hyperscalers --- Amazon Web Services, Microsoft Azure, and Google Cloud --- have moved beyond infrastructure provision into managed AI services, industry-specific solutions, and orchestration capabilities \citep{Jacobides2022}. Each layer they add absorbs work that specialist firms previously differentiated on. Large systems integrators --- Accenture, TCS, Capgemini --- commoditize delivery through scale, acquiring dozens of specialist firms per year and building proprietary platforms that standardize what boutique firms once offered as craft \citep{Gawer2022DigitalAge}. Software-as-a-Service vendors --- Oracle, SAP, ServiceNow --- close the application gap with autonomous, AI-native products where clients subscribe directly, bypassing the integrator entirely and converting what were custom implementation projects into configuration exercises.

Most recently, generative AI --- ChatGPT, Claude, Copilot --- commoditizes knowledge work itself: the highest-margin activity in IT services. Strategy advice, architecture reviews, security assessments, and code generation --- activities that specialist firms charged premium rates for because they required human expertise --- are becoming available at near-zero marginal cost. The knowledge asymmetry that justified the consulting premium is collapsing. Software development itself is being compressed: agentic coding tools mean fewer engineers deliver equivalent output, and every client will ask why they need a fifty-person team when ten people with AI agents produce the same results.

For firms operating without platform ownership --- those who cannot compete on infrastructure scale, talent volume, product IP, or generic knowledge --- the strategic literature offers a stark diagnosis. \citet{Porter1985} identifies this as ``stuck in the middle'': firms that achieve neither cost leadership nor differentiation face systematic underperformance. Transaction cost economics \citep{Williamson1991} predicts that such firms will behave opportunistically to survive, maximizing extraction from each client relationship. Ecosystem strategy \citep{Gawer2014, Jacobides2018} theorizes only two viable roles: platform owner or complementor. The non-platform orchestrator is not a recognized category.

Yet empirically, some firms do the opposite of what theory predicts. They invest in building client capabilities beyond contractual scope. They coordinate competing vendors in the client's interest. They resolve organizational paradoxes that the client cannot resolve internally. They orchestrate complex multi-vendor ecosystems without owning the platform. They act as stewards --- absorbing risk and building client independence rather than dependency.

\subsection{The Theoretical Gap}

This observed behavior violates five core theoretical assumptions simultaneously.

First, the opportunism assumption in transaction cost economics \citep{Williamson1991, Williamson1993OpportunismCritics}. TCE assumes suppliers minimize effort and maximize extraction within contractual boundaries. The observed behavior --- investing in client capabilities at the supplier's short-term cost --- contradicts this foundational behavioral assumption. This is not a parameter adjustment; relaxing the opportunism assumption fundamentally changes the theory.

Second, the internal control assumption in dynamic capabilities theory \citep{Teece1997, Teece2018}. Dynamic capabilities are theorized as internal organizational processes for sensing, seizing, and reconfiguring resources. No framework exists for capabilities that operate \emph{across} organizational boundaries on behalf of others --- what we term meta-dynamic capabilities. The supplier builds capabilities inside the client organization, using knowledge accumulated across its portfolio, through a mechanism that existing DC theory cannot describe.

Third, the platform orchestration assumption in ecosystem strategy \citep{Gawer2014, Adner2017, Jacobides2018}. Ecosystem theory recognizes platform owners who orchestrate through architectural control and complementors who participate. The non-platform orchestrator --- a firm that coordinates ecosystem participants through relational governance and contextual knowledge rather than platform ownership --- is not a theoretical category.

Fourth, the intra-firm boundary of stewardship theory \citep{Davis1997TowardManagement, Donaldson1991}. Stewardship theory extends agency theory by explaining why managers act as stewards of organizational interests rather than self-interested agents. This extension has not been applied to inter-organizational relationships where the steward is an external supplier acting in the client's interest without hierarchical authority or equity ownership.

Fifth, the internal resolution assumption in paradox theory \citep{Smith2011}. Organizational paradoxes --- tensions between innovation and stability, exploration and exploitation, change and continuity --- are theorized as resolvable through internal managerial capabilities. No framework exists for external resolution agents who navigate these tensions from outside, enabled by political neutrality and cross-client pattern recognition.

No single existing construct captures the combination of stewardship motive, external capability building, non-platform orchestration, cross-boundary dynamic capabilities, and external paradox resolution. Adjacent constructs address fragments: the relational view \citep{DyerSingh1998} explains inter-firm competitive advantage through relationships but assumes symmetry; boundary spanning captures cross-boundary activity but not the stewardship motive; coopetition \citep{nalebuff1996coopetition} addresses simultaneous competition and cooperation but not the asymmetric steward role.

\subsection{Contextual Intimacy as the Defensible Position}

If the four compression forces eliminate infrastructure, delivery, applications, and generic knowledge as sources of competitive advantage, what remains? We argue that the only defensible position for the non-platform supplier is \emph{contextual intimacy}: deep, trust-based understanding of how a specific client needs to evolve within its specific ecosystem --- given its politics, its legacy systems, its risk appetite, and its competitive dynamics.

This contextual knowledge has four properties that make it resistant to commoditization. It is too embedded in the client's organizational reality to transfer to a hyperscaler. It requires years of relationship to accumulate, making it unreplicable by a large integrator on a project timeline. It requires judgment across organizational boundaries that no product vendor can automate. And it depends on access to implicit organizational knowledge --- the politics, the trust relationships, the history of failed initiatives, the unspoken constraints --- that does not exist in any training corpus and cannot be reproduced by generative AI.

Contextual intimacy is not a static asset. It compounds: every interaction deepens understanding, every stewardship act increases trust, increased trust increases access, and increased access increases knowledge. This flywheel is the economic logic of stewardship --- and it is precisely the flywheel that extraction destroys.

\subsection{Research Question and Contribution}

This paper addresses the theoretical gap through abductive theory building \citep{DuboisGadde2002}: starting from the empirical anomaly, systematically demonstrating theoretical inadequacy through a structured literature review, and proposing a new construct. Our research question:

\begin{quote}
\emph{How do IT suppliers orchestrate innovation ecosystems without platform ownership, and what capabilities and governance mechanisms enable this role?}
\end{quote}

We introduce \textbf{Ecosystem Stewardship} as a theoretical construct defined by four mechanisms: External Paradox Resolution, Portfolio Transfer, Boundary Management, and Innovation Risk Orchestration. Each mechanism extends an existing theoretical stream by relaxing a specific assumption that blocks explanation of the observed behavior. The construct is developed following \citeauthor{Whetten1989}'s (\citeyear{Whetten1989}) criteria for theoretical contribution and grounded in a systematic literature review of [N] papers across five theoretical streams.

The contribution is threefold. First, we challenge the assumption --- shared across TCE, DC theory, and ecosystem strategy --- that effective orchestration requires either platform control or contractual authority. Second, we extend stewardship theory from intra-firm to inter-organizational relationships, providing a theoretical foundation for the empirically observed but theoretically unexplained behavior of external actors who steward client ecosystems. Third, we provide boundary conditions specifying when ecosystem stewardship creates value (complex ecosystems, capability gaps, trust-based relationships) and when it fails (commoditized services, power asymmetry, transactional relationships).

The paper proceeds as follows. Section II describes our abductive methodology and systematic literature review protocol. Section III synthesizes five literature streams, demonstrating the theoretical inadequacy of each. Section IV presents the Ecosystem Stewardship construct with its four mechanisms and testable propositions. Section V discusses implications, boundary conditions, and limitations. Section VI concludes with an empirical agenda for validation.

\section{Method: Abductive Theory Building}

\subsection{Epistemological Position}

This paper does not follow the positivist meta-review tradition where a systematic literature review synthesizes existing findings to derive cumulative knowledge. A standard SLR answers ``what does the literature say?'' --- it cannot produce a new theoretical construct, because if the construct were already in the literature, it would not be new.

Instead, this paper uses \textbf{abductive reasoning} \citep{DuboisGadde2002}: starting from a surprising empirical observation (IT suppliers orchestrating client ecosystems without platform ownership), demonstrating that no existing theory can explain it (through a systematic literature review), and proposing a new construct that, if true, would make the observed behavior intelligible.

The SLR serves as evidence of theoretical inadequacy --- it demonstrates the gap --- not as the generative method. The theoretical contribution comes from the abductive inference, not from summarizing what others wrote.

\subsection{Research Design}

The research proceeds through five phases:

\begin{enumerate}
    \item \textbf{Phase 0: Empirical puzzle.} Formal articulation of the anomaly --- five observed supplier behaviors that violate core assumptions of existing theory (Section I).
    \item \textbf{Phase 1: Systematic search.} Three query clusters across Scopus, Web of Science, and IEEE Xplore, targeting five theoretical streams.
    \item \textbf{Phase 2: Screening and snowballing.} Deduplication, keyword screening, and citation expansion (backward and forward) until saturation.
    \item \textbf{Phase 3: Theoretical coding.} Dual-pass extraction coding each paper's theoretical assumptions, supplier role characterization, and anomaly coverage.
    \item \textbf{Phase 4: Abductive synthesis.} Concept matrix construction, gap visualization, construct definition, and proposition formulation.
\end{enumerate}

This design combines elements from four established methodological traditions: concept-centric literature review \citep{WebsterWatson2002}, cross-case display \citep{MilesHuberman1994}, framework synthesis \citep{Carroll2013}, and abductive case research \citep{DuboisGadde2002}. Each component is well-established; their combination in service of abductive theory building is the methodological contribution.

\subsection{Empirical Anchoring}

The empirical puzzle is anchored in practitioner observation at Schuberg Philis (SBP), a Dutch IT services firm. The first author serves as CIO at SBP, making this insider practitioner research \citep{BrannickCoghlan2007}. This position creates unique access to the phenomenon --- the behaviors in A1--A5 are directly observed, not inferred from secondary data --- but also requires explicit reflexivity about potential bias.

Three safeguards address this: (1) the SLR is conducted independently of the motivating case --- all literature coding follows a pre-defined protocol; (2) the construct definition follows \citeauthor{Suddaby2010}'s (\citeyear{Suddaby2010}) criteria, which require discriminant validity against existing constructs; (3) the propositions are formulated to be testable through empirical validation in Paper 3, which will survey multiple firms beyond SBP.

\subsection{Search Strategy}

\subsubsection{Query clusters}

Three search strings target the theoretical pillars of the argument:

\begin{itemize}
    \item \textbf{Query A --- Supplier role evolution:} \textit{(``IT supplier'' OR ``IT vendor'' OR ``service provider'' OR ``IT partner'' OR ``outsourcing provider'') AND (``ecosystem'' OR ``orchestrat*'' OR ``steward*'' OR ``co-creation'' OR ``value creation'')}
    \item \textbf{Query B --- Dynamic capabilities beyond boundaries:} \textit{(``dynamic capabilities'' OR ``meta-capabilities'' OR ``higher-order capabilities'') AND (``supplier'' OR ``vendor'' OR ``external'' OR ``boundary-spanning'' OR ``inter-organizational'')}
    \item \textbf{Query C --- Ecosystem governance:} \textit{(``ecosystem'' OR ``platform ecosystem'' OR ``innovation ecosystem'') AND (``orchestrat*'' OR ``governance'' OR ``coordination'' OR ``complementor'') AND (``supplier'' OR ``vendor'' OR ``partner'')}
\end{itemize}

\subsubsection{Databases and scope}

Searches were conducted across Scopus (broadest coverage of management and IS literature), Web of Science (high-impact journals), and IEEE Xplore (IT-specific ecosystem and platform papers). The date range spans 2005--2026, anchored by \citet{Teece1997}'s foundational dynamic capabilities work while including recent ecosystem strategy developments.

\subsubsection{Inclusion and exclusion criteria}

Papers were included if they are peer-reviewed journal articles or top-tier conference papers, published in English, addressing supplier or vendor roles in innovation, ecosystems, or capability building. Papers were excluded if they focus on pure consumer/B2C contexts, domain-specific applications (medical, agriculture, education, military), or lack theoretical contribution (pure case description without theoretical framing).

\subsection{Screening Protocol}

\subsubsection{Deduplication}

Papers appearing in multiple databases were deduplicated using three methods in priority order: DOI exact match, title-plus-author fuzzy match (threshold 0.90), and title-only fuzzy match (threshold 0.95).

\subsubsection{Automated screening}

Title and abstract screening used keyword-based classification against four inclusion categories (ecosystem, supplier role, capabilities, governance) and domain-based exclusion lists. Papers passing the keyword gate were ranked by semantic similarity to the research question using embedding-based classification.

\subsubsection{Screening reliability}

To validate automated screening, a random sample of 50 papers was independently classified by the first author. Agreement was measured using Cohen's Kappa ($\kappa$). A threshold of $\kappa > 0.80$ (almost perfect agreement) was required to proceed; $\kappa$ between 0.60 and 0.80 triggered prompt revision and re-screening.

\subsubsection{Citation expansion}

Backward snowballing extracted reference lists from included papers via CrossRef, filtering for journal articles and iterating until saturation ($<$ 5\% new papers per iteration, maximum three iterations). Forward snowballing identified papers citing each included paper via OpenAlex, limited to one iteration.

\subsection{Theoretical Coding}

\subsubsection{Dual-pass extraction}

Each included paper was processed through two extraction passes. Pass 1 captured standard metadata, research question, methodology, key findings, and empirical patterns using the CAMO framework (Context-Agency-Mechanism-Outcome). Pass 2 --- the primary coding for this paper --- captured the theoretical assumptions underlying each paper's treatment of suppliers.

\subsubsection{Coding scheme}

The theoretical coding scheme was built from the anomaly table in Phase 0. For each paper, six dimensions were coded:

\begin{enumerate}
    \item \textbf{Theoretical lens:} primary and secondary streams (TCE, dynamic capabilities, ecosystem strategy, stewardship theory, paradox theory, relational view, other).
    \item \textbf{Supplier role:} how the paper characterizes the supplier (vendor, partner, orchestrator, steward, platform owner, complementor, not discussed).
    \item \textbf{Supplier motive assumed:} the paper's implicit or explicit assumption about supplier motivation (opportunistic, self-interested, neutral, collaborative, benevolent).
    \item \textbf{Capability scope:} whether capabilities discussed are internal only, inter-organizational, or meta-capabilities operating on behalf of others.
    \item \textbf{Orchestration mode:} the coordination mechanism assumed (platform-based, contractual, relational, hierarchical, hybrid).
    \item \textbf{Anomaly coverage:} for each of the five anomalies (A1--A5), whether the paper addresses, describes, contradicts, or is silent on the behavior. An anomaly score (0--5) was computed as the count of ``addresses'' and ``describes'' classifications.
\end{enumerate}

Additionally, each paper was coded for its \textit{blocking assumption} --- the specific theoretical assumption that prevents it from explaining ecosystem stewardship --- and its \textit{nearest concept} --- the construct that comes closest to ecosystem stewardship without capturing it fully.

\subsection{Analysis: The Concept Matrix}

The analysis centers on a \textbf{concept matrix} \citep{WebsterWatson2002} that cross-tabulates supplier role (rows) against theoretical lens (columns), following Miles and Huberman's (\citeyear{MilesHuberman1994}) cross-case display logic. This standard method for concept-centric literature reviews in information systems is extended with two features from the abductive methodology.

\textbf{Extension 1: Anomaly score overlay.} Each cell is annotated with the average anomaly score of papers it contains. This follows the framework synthesis approach of \citet{Carroll2013}, where data is coded against a pre-defined framework --- in this case, the empirical puzzle from Phase 0. Three patterns emerge:

\begin{itemize}
    \item \textit{Low anomaly score in dense cells:} well-theorized territory --- the theory adequately explains what is observed.
    \item \textit{High anomaly score in dense cells:} \textbf{misclassification} --- the phenomenon is observed but forced into an inadequate theoretical lens. These are the strongest evidence for the gap.
    \item \textit{Empty cells:} under-explored territory requiring verification (did the search miss relevant papers, or does the gap genuinely exist?).
\end{itemize}

Misclassified cells --- papers that describe steward-like behavior but code it as ``partnership'' under TCE because the theory has no steward category --- are more valuable as gap evidence than empty cells. They demonstrate that the phenomenon exists \textit{and} that existing theory forces researchers to mislabel it.

\textbf{Extension 2: Blocking assumption annotation.} For each theoretical stream, the matrix is annotated with the specific assumption that prevents papers in that stream from explaining ecosystem stewardship. This follows Gioia et al.'s (\citeyear{Gioia2013}) first-order to second-order coding structure: the paper's own characterization (first-order) is mapped to the blocking assumption (second-order) and then to the aggregate theoretical dimension it challenges.

\subsection{Construct Development}

The construct definition follows \citeauthor{Suddaby2010}'s (\citeyear{Suddaby2010}) construct clarity criteria: (1) what the construct is, (2) what it is not (discriminant validity against adjacent constructs), (3) boundary conditions specifying when it applies, and (4) its relationship to existing constructs. Propositions were formulated to be traceable to at least three papers from the SLR, with counter-evidence explicitly addressed for each proposition. The framework targets \citeauthor{Whetten1989}'s (\citeyear{Whetten1989}) criteria for theoretical contribution: what (the construct and mechanisms), how (the causal logic), why (the structural argument), and who/where/when (the boundary conditions).

\section{Literature: Five Streams, One Gap}

% 3a: Transaction Cost Economics — Williamson 1985, 1991, 1993
%     Supplier as opportunistic actor. Fails: assumes value extraction.
% 3b: Dynamic Capabilities — Teece 1997, 2007
%     Internal control. Fails: no cross-boundary capability operation.
% 3c: Ecosystem Strategy — Adner 2017, Jacobides 2018, Gawer 2014
%     Platform owners and complementors. Fails: non-platform orchestrators not a category.
% 3d: Stewardship Theory — Davis et al. 1997, Hernandez 2012
%     Intra-firm stewardship. Fails: not extended to inter-organizational.
% 3e: Paradox Theory — Smith & Lewis 2011
%     Internal resolution. Fails: no external resolution agents.
% 3f: Industry Architecture — Jacobides, Knudsen & Augier 2006
%     Value migration across layers. The structural context.

% Thematic matrix: supplier role × theoretical lens
% The gap is at the intersection — no construct covers all 5 anomaly dimensions

\section{Ecosystem Stewardship: Construct and Mechanisms}

% Construct definition following Suddaby (2010):
%   What it is / What it is not / Boundary conditions / Relation to existing constructs

% Mechanism 1: External Paradox Resolution (extends Smith & Lewis 2011)
% Mechanism 2: Portfolio Transfer (extends Teece 2007)
% Mechanism 3: Boundary Management (extends Tushman 1977)
% Mechanism 4: Innovation Risk Orchestration (extends Jacobides 2018)

% Propositions: testable claims, each traceable to >= 3 SLR papers
% Counter-evidence addressed per proposition

\section{Discussion}

% What this framework explains that existing theory cannot
% The structural argument: when does stewardship emerge?
% Boundary conditions: when stewardship fails
% Implications for TCE, DC, ecosystem theory, stewardship theory, paradox theory
% RBV implications (secondary)
% Reflexivity: insider practitioner research — strengths and limitations

\section{Conclusion}

% Contribution summary
% Limitations (single-firm empirical anchor, abductive method)
% Empirical agenda -> Paper 3 (survey + cases across multiple firms)


%----------------------------------------------------------------------------------------
%	BIBLIOGRAPHY
%----------------------------------------------------------------------------------------
\printbibliography[title=References]

%----------------------------------------------------------------------------------------
\end{document}
